\SourceLeaf{446}%
\begin{quote}
»If shortest lines of equal length are drawn from a point on the curved surface, they meet the line joining their endpoints everywhere at right angles«.
\end{quote}
\noindent\hspace*{1.2em}One can prove in much the same way that if any curved line is given on the curved surface and shortest lines of equal length are drawn from each of its points on the same side of the given curve and perpendicular to it, the line joining their endpoints is cut by the shortest lines everywhere at right angles. One need only let $\theta$ in the previous derivation denote the length along the given curved line from an arbitrary point. The previous calculations then remain valid, except that the value $u=0$ for $s=0$ now follows from the hypothesis itself.

\Needspace{5\baselineskip}
\begin{center}18.\end{center}
\noindent\hspace*{1.2em}The relations between the quantities in these constructions deserve to be worked out in more detail. First, writing $m$ for $\frac{\partial t}{\partial\theta}$ for brevity, we have
\begin{equation}
\frac{\partial x}{\partial s}=\xi,\qquad\frac{\partial y}{\partial s}=\eta,\qquad\frac{\partial z}{\partial s}=\zeta,\tag{1}
\end{equation}
\begin{equation}
\frac{\partial x}{\partial\theta}=m\xi',\qquad\frac{\partial y}{\partial\theta}=m\eta',\qquad\frac{\partial z}{\partial\theta}=m\zeta',\tag{2}
\end{equation}
\begin{equation}\xi\xi+\eta\eta+\zeta\zeta=1\tag{3}\end{equation}
\begin{equation}\xi'\xi'+\eta'\eta'+\zeta'\zeta'=1\tag{4}\end{equation}
\begin{equation}\xi\xi'+\eta\eta'+\zeta\zeta'=0.\tag{5}\end{equation}
Furthermore,
\begin{equation}XX+YY+ZZ=1\tag{6}\end{equation}
\begin{equation}X\xi+Y\eta+Z\zeta=0\tag{7}\end{equation}
\begin{equation}X\xi'+Y\eta'+Z\zeta'=0\tag{8}\end{equation}
and
\begin{equation}
\left\{\begin{aligned}X&=\zeta\eta'-\eta\zeta'\\Y&=\xi\zeta'-\zeta\xi'\\Z&=\eta\xi'-\xi\eta'\end{aligned}\right.\tag*{[9]}
\end{equation}
